\documentclass{situationist-preprint}

\title{Progress in Art}
\msid{AES-16-a}
\msfamily{AES-16}
\msdiscipline{Aesthetics}
\mstype{Principal manuscript}
\mskeywords{progress; art history; artistic means; artistic value; Vasari; Gombrich; Danto; achievement}

\begin{document}
\maketitle

\begin{abstract}
Is there progress in art? Since Vasari, histories of art have described the acquisition of perspective, anatomy and naturalistic light as an advance, and in the sciences later work is in general better than earlier work in the respects that matter. We argue that art differs: there is progress in artistic means but not in artistic value. Artistic means, such as techniques, materials, instruments and representational devices, accumulate, and later artists command means that earlier ones lacked. But the artistic value of a work consists in its aesthetic value, appreciated in its correct category, and in the value of the achievement it embodies, assessed relative to its maker's situation. Neither component increases with a work's position in history: the new means of a later period define new categories rather than raising the standard of the old, and achievement is measured against the means available. We show that the progress narratives of Vasari and Gombrich are narratives of means, that Danto's narrative of art's self-understanding is not a narrative of value, and that a growing stock of great works is accumulation, not progress.
\end{abstract}

\section{The problem}

\citet{Vasari1568} organized the history of Italian art as a progress in three ages, from the recovery of design by Giotto, through the conquest of perspective and anatomy in the fifteenth century, to the perfection achieved by Leonardo, Raphael and Michelangelo. The narrative served a view of art as the imitation of nature, in which success is measurable and cumulative. The Quarrel of the Ancients and the Moderns, at the end of the seventeenth century, turned on the same question for literature: Perrault argued, in his \emph{Parallèle des Anciens et des Modernes}, that the moderns, having more knowledge, surpassed the ancients, and his opponents denied it. \citet{Hume1757} offered a test that told against Perrault: the same Homer who pleased at Athens and Rome two thousand years ago is still admired, whereas the systems of the philosophers of those ages have been overturned.

In the twentieth century, \citet{Gombrich1960} gave the progress narrative a sophisticated form: the history of representation is a history of \enquote{making and matching}, in which artists modify inherited schemata in the light of their effects, and the naturalistic achievements of Greek and Renaissance art are genuine discoveries. Gombrich was careful to say that these discoveries are discoveries of means, and that a more naturalistic painting is not for that reason a better one. \citet{Danto1984,Danto1997} proposed a different narrative, on which the history of Western art, after the abandonment of the mimetic project, became a progress toward art's understanding of its own nature, which was completed when works such as Warhol's \emph{Brillo Box} showed that art and non-art can be perceptually indiscernible. On Danto's view, that history has ended, and art now proceeds without a narrative of progress.

We argue that there is progress in artistic means and no progress in artistic value. The contrast with morality is instructive: MTE-19 argues that moral progress is real, as the expansion of the moral circle. We show why nothing corresponding holds for the value of art.

\section{Framework}

\begin{definition}[Progress]\label{d:prog}
There is \emph{progress} in a domain with respect to a dimension $D$ if, as a rule, later states of the domain stand higher on $D$ than earlier states, and they do so in virtue of what was achieved in the earlier states.
\end{definition}

The second clause distinguishes progress from mere improvement and from accident: progress is cumulative, with later states building on earlier ones.

\begin{definition}[Artistic means]\label{d:means}
The \emph{artistic means} of a period are the techniques, materials, instruments, forms and representational devices available to its artists: linear perspective, oil glazing, the piano, equal temperament, the novel's free indirect style, synchronized sound in film.
\end{definition}

\begin{principle}[Artistic value]\label{p:value}
The artistic value of a work consists in its aesthetic value, as perceived by qualified observers in its correct category, together with the value of the achievement it embodies, assessed relative to its maker's situation, including the means then available (AES-02, AES-04).
\end{principle}

\begin{principle}[Categories of means]\label{p:catmeans}
New artistic means create new categories in which works are perceived, rather than raising the standard by which works in existing categories are perceived.
\end{principle}

Principle~\ref{p:value} combines the response-dependent account of aesthetic value in AES-02, on which qualified observers perceive works in their relevant categories, with the account of artistic value in AES-04, on which the achievement a work embodies is relative to its maker's situation. Principle~\ref{p:catmeans} applies Walton's thesis that aesthetic properties depend on the category in which a work is correctly perceived \citep{Walton1970}. The flatness of a medieval panel is standard for its category and contributes to its hieratic character; it is not a failed attempt at perspective, and the invention of perspective did not make it worse.

\section{The argument}

\begin{lemma}\label{l:means}
There is progress in artistic means.
\end{lemma}

\begin{proof}
Artistic means accumulate. A technique, once invented, is transmitted through workshops, treatises, academies and the works themselves, and later artists can use it or refuse it. Later periods therefore command, as a rule, a larger and more powerful repertoire of means, in the sense that they can achieve every effect earlier means achieved, by using those means, and further effects besides; and they command it in virtue of what earlier artists invented. By Definition~\ref{d:prog}, this is progress with respect to the range of means.
\end{proof}

Losses occur: techniques are sometimes forgotten, as the recipes of certain medieval pigments were. They are exceptions to a rule, which Definition~\ref{d:prog} allows.

\begin{lemma}\label{l:aes}
There is no progress in aesthetic value.
\end{lemma}

\begin{proof}
By Principle~\ref{p:value}, a work's aesthetic value is fixed by qualified observers perceiving it in its correct category. By Principle~\ref{p:catmeans}, the new means of later periods create new categories and do not change the standard of perception in the categories of earlier works. A later work therefore does not stand higher in aesthetic value than an earlier one in virtue of its later means: the earlier work is perceived in a category in which those means are not standard, and its aesthetic value is unaffected by their existence. Since the means are what later periods inherit, there is no dimension of aesthetic value on which later works stand higher in virtue of what earlier works achieved, and by Definition~\ref{d:prog} no progress.
\end{proof}

\begin{lemma}\label{l:ach}
There is no progress in achievement value.
\end{lemma}

\begin{proof}
By Principle~\ref{p:value}, achievement is assessed relative to the maker's situation, including the means available. A later artist who uses means invented earlier achieves nothing in virtue of using them that the inventor did not achieve in inventing them; and an artist who invents new means achieves something relative to the means she inherited, as her predecessors did relative to theirs. Inheriting more means raises the baseline against which achievement is measured as much as it raises what can be done. So the achievement embodied in works does not, as a rule, increase with their position in history.
\end{proof}

\begin{theorem}\label{t:main}
There is progress in artistic means but not in artistic value.
\end{theorem}

\begin{proof}
The first conjunct is Lemma~\ref{l:means}. By Principle~\ref{p:value}, artistic value consists of aesthetic value and achievement value, and by Lemmas~\ref{l:aes} and~\ref{l:ach}, neither progresses. So artistic value does not progress.
\end{proof}

Theorem~\ref{t:main} explains Hume's observation. Homer is still admired because the value of the \emph{Iliad} is the value of a work perceived in its category and assessed as an achievement in its situation, and neither is altered by the later invention of the novel. The overturning of ancient science, by contrast, is the overturning of claims judged by a fixed aim, truth about the world, against which later means do yield better results. The same contrast separates art from morality. Moral progress, on the account of MTE-19, is measured against a fixed standard, the scope of the moral circle. Art has no fixed aim of this kind: its aims are set in part by its categories, and new means bring new categories with new aims.

\section{Objections and replies}

\begin{objection}[Mimesis]\label{o:mimesis}
Representational accuracy is an aim of much art, and later art achieves it better. Vasari's narrative is a narrative of value \citep{Vasari1568}.
\end{objection}

\begin{reply}
Representational accuracy is a means to the aims of representational art, not the measure of its value. If it were, a photograph would surpass every painting, and AES-14 argues that photography is art through framing and selection, not through accuracy. \citet{Gombrich1960} drew the right conclusion from his own history: the discoveries of naturalism are discoveries of means, and the value of a Giotto is not diminished by Masaccio's perspective.
\end{reply}

\begin{objection}[Progress within traditions]\label{o:local}
Within a tradition, later works solve problems earlier works posed and could not solve: the development of sonata form from Haydn to Beethoven is a sequence of solutions. This is progress in value.
\end{objection}

\begin{reply}
Problem-solving within a tradition is the invention of means relative to aims the tradition sets. The later solution has achievement value relative to its situation, and so did the earlier formulation of the problem. Beethoven's later solutions do not make Haydn's symphonies worse, nor are they better in virtue of coming later; they are works with a different situation and, often, a different category. Local sequences of solutions are progress in means.
\end{reply}

\begin{objection}[Danto]\label{o:danto}
The history of art has a direction: toward the philosophical self-understanding of art. This is progress \citep{Danto1984}.
\end{objection}

\begin{reply}
If Danto is right, the history he describes is progress in understanding what art is, which is a cognitive achievement. It is not progress in artistic value, and Danto did not claim it was: \emph{Brillo Box} is not a better work than a Rembrandt in virtue of its place in the narrative. His narrative is consistent with Theorem~\ref{t:main} and indeed requires its distinction, since its endpoint is a pluralism in which no style is historically mandated and none is superior for coming later.
\end{reply}

\begin{objection}[The stock]\label{o:stock}
The number of great works increases over time, and each period inherits all earlier ones. The world of art is therefore richer in value than it was.
\end{objection}

\begin{reply}
The stock grows, but this is accumulation, not progress in artistic value as Definition~\ref{d:prog} defines it. Later works do not stand higher than earlier ones; there are simply more works. The growth of the stock is, moreover, offset by losses of works that do not survive, and in any case it concerns the totality of art and not the value of the works that make it up.
\end{reply}

\section{Conclusion}

There is progress in art, but it is progress in means: techniques, materials, instruments and forms accumulate, and later artists command more of them. There is no progress in artistic value. Works are valued as perceived in their categories and as achievements relative to their situations, and new means bring new categories and new baselines rather than higher standards for the old. That is why Homer still pleases, why naturalistic discovery did not demote Giotto, and why the history of art, unlike the history of science or of morality, has no direction in value.

\begin{thebibliography}{99}\small
\bibitem[Danto(1984)]{Danto1984} Danto, A.~C. (1984). The end of art. In B. Lang (Ed.), \emph{The Death of Art}. Haven Publishers.
\bibitem[Danto(1997)]{Danto1997} Danto, A.~C. (1997). \emph{After the End of Art: Contemporary Art and the Pale of History}. Princeton University Press.
\bibitem[Gombrich(1960)]{Gombrich1960} Gombrich, E.~H. (1960). \emph{Art and Illusion: A Study in the Psychology of Pictorial Representation}. Pantheon Books.
\bibitem[Hume(1757)]{Hume1757} Hume, D. (1757). Of the standard of taste. In \emph{Four Dissertations}. A. Millar.
\bibitem[Vasari(1568)]{Vasari1568} Vasari, G. (1568). \emph{Le vite de' più eccellenti pittori, scultori e architettori}. 2nd ed. Giunti.
\bibitem[Walton(1970)]{Walton1970} Walton, K.~L. (1970). Categories of art. \emph{The Philosophical Review}, 79(3), 334--367.
\end{thebibliography}

\end{document}
